\section{Correctness Proofs}
\label{sec:appendix-correctness}
The following proof sketches address the protocol safety properties stated in the paper under its fixed-membership, crash-recovery assumptions. The accompanying TLA$^+$ artifact provides machine-checked proofs and supplementary finite model checks.

\noindent\textit{Precise recovery semantics.}
The recovery rule in Section~\ref{subsubsec:traft-recovery} is interpreted
as follows. A candidate collects reports from a majority, including itself,
after the reporters adopt its term. Each report contains the replica's
installed log and its complete durable recovery proposals, including those
whose confirmation is unknown. For a proposal $p$, let $B(p)$ be its
retained prefix together with its opening no-op, and let $t_p$ and
$\rho_p$ be that no-op's creation term and partition index. A proposal is
retained only after all of $B(p)$ is durable. Among reported proposals,
select one with lexicographically maximal $(t_p,\rho_p)$. Let $U$ be the
union of all installed log entries and all complete recovery results
contained in the collected reports. The recovered log is
\[
  L = B(p) \cup
      \{e\in U : \operatorname{term}(e)=t_p
                     \land \operatorname{part}(e)\geq\rho_p\}.
\]
Thus, $B(p)$ is inherited in full, and reports from replicas at different
ordinary-partition progress contribute extensions from its creation term.
Original entry identities, indices, creation terms, partition indices,
and contents are preserved. If no recovery proposal is reported, recovery
starts from an empty log only when the reports also contain no entries.

The candidate appends a new-term no-op immediately after $L$ in the next
partition. Replicas may copy the result incrementally, but retain and
acknowledge its recovery proposal only after the complete result is
durable. Acknowledgment does not require prior commitment. After a majority
has acknowledged, the candidate installs the result, commits its entries
in partition order, and serves writes only after the recovered prefix and
the new no-op are committed. See
\texttt{Latest}, \texttt{ReportedEntries}, \texttt{Recovered}, and
\texttt{BuildRecovery}, together with the recovery actions in
\href{https://github.com/LittleHealth/traft/blob/traft-proof-v1/tlaplus-proof/models/TRaftCore.tla}{\texttt{TRaftCore.tla}}
(tag \texttt{traft-proof-v1}).

\subsection{Log Matching (Property~\ref{prop:log-matching})}
\label{subsec:proof-log-matching}

\begin{proof}[Proof sketch]
Before serving in a new term, quorum-confirmed recovery establishes a
common prefix and a new-term no-op. Within that term, the leader assigns
fresh entry indices and opens ordinary successor partitions, fixing each
preceding partition's membership. Under this shared recovery origin, a
follower's entries in the preceding partition are a subset of the leader's
fixed contents. Equality of its local distinct-entry count with the
leader's final total and committed counts therefore establishes both
local completeness and full commitment of that partition. The follower
then persists the successor entry, advances its current partition, and
acknowledges it; an ordinary boundary requires no separate recovery-proposal
confirmation round. Commitment of the new entry still requires a majority
of durable acknowledgments in the current term and commitment of all
preceding partitions. Subsequent replication preserves the closed prefix.
Across terms, the recovery-preservation argument in the proof of
Property~\ref{prop:leader-completeness} retains that prefix,
including through interrupted recovery. Thus matching terms and partition
indices identify identical preceding partitions, while copies of the
open partition may differ.
\end{proof}

\subsection{Leader Completeness (Property~\ref{prop:leader-completeness})}
\label{subsec:proof-leader-completeness}

\begin{proof}[Proof sketch]
First consider the complete result of a quorum-acknowledged recovery proposal.
Every later recovery quorum intersects its acknowledgment quorum. Term
fencing ensures that the shared replica durably retained the proposal before
reporting the higher term, and retained proposals are reported even if
confirmation is unknown. The selected proposal is this proposal or a
later one: the protocol creates at most one recovery proposal per term,
and proposals from intervening terms preserve it by induction. Ordinary
partition advances extend the installed log without creating new recovery
proposals. This argument covers all
recovery attempts, including interrupted ones, and preserves closed
contents even before a partition is committed as a whole.

For a committed entry $e$ not yet covered by such a proposal, suppose the
first later recovery omits it. A replica in the intersection of the
recovery and acknowledgment quorums reports the durable history containing
$e$, together with its retained recovery proposals. Its recovery origin
can precede $e$'s ordinary partition; an ordinary boundary creates no
separate proposal.
Recovery either inherits $e$ in the selected recovery result or includes it
in the reported extensions defined above. Selecting a later-term proposal cannot
omit $e$ without an earlier violating recovery, contradicting the choice
of the first omission. Every valid recovery result therefore preserves
committed entries, including those in open partitions. A Leader confirms
this result before serving, and Leader Append-Only preserves it thereafter.
Together, these arguments establish both partition and entry preservation.
\end{proof}

\subsection{State Machine Safety (Property~\ref{prop:state-machine-safety})}
\label{subsec:proof-state-machine-safety}

\begin{proof}[Proof sketch]
Suppose two replicas apply different entries at the same log index.
Both entries must be committed. In the same authorizing term, the
LeaderCandidate's fixed recovery result and fresh-index assignment exclude
two different entries at that index. In different terms, the preservation
argument for Property~\ref{prop:leader-completeness} retains
the earlier committed entry in every later valid recovery result,
including interrupted attempts; replication cannot replace it.
Both cases give a contradiction. The fixed-prefix preservation
prevents two fully applied partitions at one partition index from having
different contents.

For execution-result agreement, replicas apply preceding partitions first
and each entry at most once. Within a partition, the assumed independence
of operations makes their order immaterial. Deterministic replicas that
have applied the same committed entries therefore reach the same state;
replicas with different applied subsets may temporarily differ.
\end{proof}



\section{Theoretical Proofs}
\label{sec:appendix-proofs}
\subsection{Proof of Theorem~\ref{thm:raft-amplification}}
\begin{proof}
The leader maintains an in-memory log cache of size $M$, which can retain at
most $M/b$ average-sized entries.
Because approximately $2\ell/\Lambda$ entries are generated
during one expected round trip, the cache becomes saturated when the
RTT-scale replication working set exceeds $M/b$.
This yields the cache-saturation threshold for $\Lambda$
\[
\Lambda^*
=
\frac{2b\ell}{M}.
\]

When $\Lambda\geq\Lambda^*$, the RTT-scale working set fits in the cache, so
all follower streams can reuse one shared disk-loaded copy, resulting in one
entry-equivalent disk cost.
When $\Lambda<\Lambda^*$, the working set exceeds the cache capacity.
The entry may then need to be loaded independently for the remaining $f-1$
follower streams, adding $f-1$ entry-equivalent disk costs.
Combining the two cases yields
Equation~\ref{eq:raft-amplification}.
\end{proof}

\subsection{Proof of Theorem~\ref{thm:traft-amplification}}
\begin{proof}
For \textsc{TRaft}, consider a partition containing $\theta\geq1$ entries.
Let
$\kappa:=\min\{\theta,w\}$
denote the boundary-window size, the maximum number of entries in the
partition that may overlap with the unfinished tail of the preceding
partition.

When $\Lambda\geq\Lambda^*$, the RTT-scale working set remains in the cache,
and all entries incur only one shared disk cost.
When $\Lambda<\Lambda^*$, at most the first $\kappa$ entries of the current
partition are forced onto the cold path by the preceding partition.
For these entries, the disk-read multiplicities are
$1+(f-1)i/w$ for
$i=w-\kappa+1,\ldots,w$; the remaining $\theta-\kappa$ entries retain unit
multiplicity.
Averaging over all $\theta$ entries gives
\begin{align*}
\psi_{\mathsf{D}}^{\mathsf{T}}
&=
1+
\mathbb{I}[\Lambda<\Lambda^*]
\frac{f-1}{w\theta}
\sum_{i=w-\kappa+1}^{w}i \\
&=
1+
(f-1)
\frac{\kappa(2w-\kappa+1)}{2w\theta}
\mathbb{I}[\Lambda<\Lambda^*],
\end{align*}
which proves Equation~\ref{eq:traft-psi} as claimed.
\end{proof}

\subsection{Proof of Theorem~\ref{thm:raft-disorder-amplification}}
\begin{proof}
For one follower, each advancement of the in-flight window exposes one head
adjacent pair, which is inverted with probability $\alpha$. Under the
window-level repair model, this inversion invalidates all $w$ in-flight entries
and produces an expected repair cost of $\alpha w$ entry-equivalent
units per entry.

The follower channels are homogeneous and independent. Summing the expected
repair cost over the $f$ followers gives $f\alpha w$. Because each
retransmitted entry incurs one disk read and one network send under our
accounting convention, the same result applies to both resources.
\end{proof}

\subsection{Proof of Theorem~\ref{thm:traft-disorder-amplification}}
\begin{proof}
Entries within the same partition may be replicated in parallel without
strict-order repair. Consequently, only the $\kappa$ tail entries of the
preceding partition can overlap with the first entry of the current partition.

For one follower, the $i$-th closest entry in the preceding partition is
overtaken by the first entry of the current partition with probability
$\alpha_i$ and invalidates a suffix of length $w-i+1$. The expected number of
retransmitted entries at one boundary is therefore given by
\[
\sum_{i=1}^{\kappa}\alpha_i(w-i+1).
\]
Multiplying by the $f$ followers and normalizing by the $\theta$ entries in the
partition gives Equation~\ref{eq:traft-phi}. Each retransmitted
entry contributes equally to disk-read and network-send costs.

Moreover, since $F_{\Delta\ell}$ is nondecreasing and
$i\Lambda\ge\Lambda$,
\[
\alpha_i
=
1-F_{\Delta\ell}(i\Lambda)
\le
1-F_{\Delta\ell}(\Lambda)
=\alpha,
\]
and therefore we obtain
\begin{align*}
\phi_{\mathsf{D}}^{\mathsf{T}}
=\phi_{\mathsf{N}}^{\mathsf{T}}
&\le \frac{f}{\theta}\sum_{i=1}^{\kappa}\alpha w
\le f\alpha w.
\end{align*}
Thus, \textsc{TRaft} never incurs greater disorder-induced amplification than Raft
under the same replication-delay distribution and in-flight window size.
\end{proof}

\subsection{Expected Resource Costs}
\label{subsec:expected-resource-cost}
Combining Definition~\ref{def:replica-amplification} with the Raft
factors gives
\begin{align}
\mathcal{C}_{\mathsf{D}}^{\mathsf{R}}
&=2+(f-1)\mathbb{I}[\Lambda<\Lambda^*]+f\alpha w,\nonumber\\
\mathcal{C}_{\mathsf{N}}^{\mathsf{R}}&=1+f+f\alpha w.
\label{eq:appendix-raft-cost}
\end{align}
For TRaft, let $m=\min\{w,\Theta\}$. From
Equation~\ref{eq:partition-size-distribution}, for every $1\leq i\leq m$,
define the truncated inverse partition-size moment $H_i$ by averaging over $\theta$
\begin{align*}
\mathbb{E}_{\theta}\!\left[
\frac{\mathbb{I}[\theta\geq i]}{\theta}
\right]
&=
\sum_{k=i}^{\Theta}\frac{\Pr(\theta=k)}{k}\\
&=
\gamma\sum_{k=i}^{\Theta-1}\frac{(1-\gamma)^{k-1}}{k}
+\frac{(1-\gamma)^{\Theta-1}}{\Theta}
=H_i.
\end{align*}
Setting $i=1$ yields
$H_1=\mathbb{E}_{\theta}[1/\theta]$.

Since $\kappa=\min\{\theta,w\}$, expand each conditional boundary term
as a finite sum over $1\leq i\leq m$ using
$\mathbb{I}[\theta\geq i]$. Taking expectations and applying
Definition~\ref{def:replica-amplification} gives
\begin{align}
\mathcal{C}_{\mathsf{D}}^{\mathsf{T}}
&=2+\frac{f-1}{w}\mathbb{I}[\Lambda<\Lambda^*]
\sum_{i=1}^{m}(w-i+1)H_i\nonumber\\
&\quad+f\sum_{i=1}^{m}\alpha_i(w-i+1)H_i,\nonumber\\
\mathcal{C}_{\mathsf{N}}^{\mathsf{T}}
&=1+f+f\sum_{i=1}^{m}\alpha_i(w-i+1)H_i.
\label{eq:appendix-traft-cost}
\end{align}

\subsection{Proof of Corollary~\ref{cor:high-frequency-raft-throughput}}
\begin{proof}
As $\Lambda\to0$, the cache-saturation indicator converges to one,
$w=\min(W,\lceil2\ell/\Lambda\rceil)$ converges to $W$, and symmetry of the
difference between two i.i.d. continuous delays gives $\alpha\to1/2$.
Substituting these limits into Equation~\ref{eq:appendix-raft-cost} gives
\begin{align*}
\lim_{\Lambda\to0}\mathcal{C}_{\mathsf{D}}^{\mathsf{R}}
=\lim_{\Lambda\to0}\mathcal{C}_{\mathsf{N}}^{\mathsf{R}}
=1+f+\frac{1}{2}fW.
\end{align*}
Applying Definition~\ref{def:resource-limited-throughput} then proves
Equation~\ref{eq:high-frequency-raft-throughput}.
\end{proof}

\subsection{Proof of Corollary~\ref{cor:fully-disordered-limit}}
\begin{proof}
When $\gamma=1$, Equation~\ref{eq:partition-size-distribution} gives
$\Pr(\theta=1)=1$. Hence, $H_1=1$, $H_i=0$ for $i\geq2$, and
$\alpha_1=\alpha$. Substitution into
the cost expressions in Appendix~\ref{subsec:expected-resource-cost} gives
\begin{align*}
\left.\mathcal{C}_{\mathsf{D}}^{\mathsf{T}}\right|_{\gamma=1}
&=2+(f-1)\mathbb{I}[\Lambda<\Lambda^*]+f\alpha w
=\mathcal{C}_{\mathsf{D}}^{\mathsf{R}},\\
\left.\mathcal{C}_{\mathsf{N}}^{\mathsf{T}}\right|_{\gamma=1}
&=1+f+f\alpha w
=\mathcal{C}_{\mathsf{N}}^{\mathsf{R}}.
\end{align*}
Applying Definition~\ref{def:resource-limited-throughput} to equal costs
proves the equality.
\end{proof}

\subsection{Proof of Corollary~\ref{cor:fully-ordered-limit}}
\begin{proof}
When $\gamma=0$, Equation~\ref{eq:partition-size-distribution} gives
$\Pr(\theta=\Theta)=1$. Thus, $H_i=1/\Theta$ for every
$1\leq i\leq m$, where $m=\min\{w,\Theta\}$. Substitution into
the cost expressions in Appendix~\ref{subsec:expected-resource-cost} gives the finite-partition costs for this case
\begin{align*}
\left.\mathcal{C}_{\mathsf{D}}^{\mathsf{T}}\right|_{\gamma=0}
&=2+\frac{(f-1)m(2w-m+1)}{2w\Theta}
\mathbb{I}[\Lambda<\Lambda^*]\\
&\quad+\frac{f}{\Theta}
\sum_{i=1}^{m}\alpha_i(w-i+1),\\
\left.\mathcal{C}_{\mathsf{N}}^{\mathsf{T}}\right|_{\gamma=0}
&=1+f+\frac{f}{\Theta}
\sum_{i=1}^{m}\alpha_i(w-i+1).
\end{align*}
For sufficiently large $\Theta>w$, $m=w$. The finite sums are then independent
of $\Theta$, and all boundary-induced terms vanish at rate
$O(1/\Theta)$ for fixed $w$. Consequently,
\begin{align*}
\lim_{\Theta\to\infty}\mathcal{C}_{\mathsf{D}}^{\mathsf{T}}
&=\mathcal{C}_{\mathsf{D}}^{\mathsf{A}}=2, &
\lim_{\Theta\to\infty}\mathcal{C}_{\mathsf{N}}^{\mathsf{T}}
&=\mathcal{C}_{\mathsf{N}}^{\mathsf{A}}=1+f.
\end{align*}
Substitution into Equations~\ref{eq:appendix-traft-cost} and
\ref{eq:resource-limited-throughput} proves
Equation~\ref{eq:fully-ordered-limit}.
\end{proof}




\input{\AppendixRoot/experiments}
